\documentclass[../../main.tex]{subfiles}

\begin{document}

\section{实变量的连续实值函数}\label{section:ch01-连续函数}

\begin{definition}[连续]\label{definition:ch01-连续}
令 $E$ 为实数集, $f:E\to\mathbf{R}$ 为定义在 $E$ 上的实值函数. 我们说 $f$ 在点 $x_{0}\in E$ 连续, 若对每个 $\varepsilon>0$, 存在 $\delta>0$ 使得对所有 $x\in E$,
\begin{align*}
|x-x_{0}|<\delta\Longrightarrow|f(x)-f(x_{0})|<\varepsilon.
\end{align*}
若 $f$ 在 $E$ 的每一点连续, 我们说 $f$ 在 $E$ 上连续.

函数 $f:E\to\mathbf{R}$ 称为 Lipschitz 函数, 若存在常数 $M\geqslant 0$ 使得对所有 $x,x'\in E$,
\begin{align*}
|f(x)-f(x')|\leqslant M|x-x'|.
\end{align*}
\end{definition}

\begin{proposition}[序列连续性]\label{proposition:ch01-序列连续}
函数 $f:E\to\mathbf{R}$ 在 $x_{0}\in E$ 连续当且仅当对每个收敛于 $x_{0}$ 的序列 $\{x_{n}\}\subseteq E$, 序列 $\{f(x_{n})\}$ 收敛于 $f(x_{0})$.
\end{proposition}

\begin{proof}
假定 $f$ 在 $x_{0}$ 连续. 令 $\{x_{n}\}\subseteq E$ 收敛于 $x_{0}$. 给定 $\varepsilon>0$, 选取 $\delta>0$ 使得 $|x-x_{0}|<\delta$ 蕴涵 $|f(x)-f(x_{0})|<\varepsilon$. 由于 $x_{n}\to x_{0}$, 存在 $N$ 使得对 $n\geqslant N$, $|x_{n}-x_{0}|<\delta$. 于是对 $n\geqslant N$, $|f(x_{n})-f(x_{0})|<\varepsilon$. 因此 $f(x_{n})\to f(x_{0})$.

反之, 假定 $f$ 在 $x_{0}$ 不连续. 则存在 $\varepsilon_{0}>0$, 使得对每个 $\delta>0$, 存在 $x\in E$ 满足 $|x-x_{0}|<\delta$ 但 $|f(x)-f(x_{0})|\geqslant\varepsilon_{0}$. 对每个自然数 $n$, 取 $\delta=1/n$, 选取 $x_{n}\in E$ 满足 $|x_{n}-x_{0}|<1/n$ 但 $|f(x_{n})-f(x_{0})|\geqslant\varepsilon_{0}$. 则 $x_{n}\to x_{0}$, 但 $f(x_{n})$ 不收敛于 $f(x_{0})$. 矛盾.
\end{proof}

\begin{proposition}[开集原象刻画]\label{proposition:ch01-原象连续}
函数 $f:\mathbf{R}\to\mathbf{R}$ 连续当且仅当对每个开集 $\mathcal{O}\subseteq\mathbf{R}$, 原象 $f^{-1}(\mathcal{O})$ 是开集.
\end{proposition}

\begin{proof}
假定 $f$ 连续. 令 $\mathcal{O}$ 为开集. 令 $x_{0}\in f^{-1}(\mathcal{O})$. 则 $f(x_{0})\in\mathcal{O}$. 由于 $\mathcal{O}$ 是开的, 存在 $\varepsilon>0$ 使得 $(f(x_{0})-\varepsilon,f(x_{0})+\varepsilon)\subseteq\mathcal{O}$. 由 $f$ 在 $x_{0}$ 连续, 存在 $\delta>0$ 使得 $|x-x_{0}|<\delta$ 蕴涵 $|f(x)-f(x_{0})|<\varepsilon$. 因此 $(x_{0}-\delta,x_{0}+\delta)\subseteq f^{-1}(\mathcal{O})$. 所以 $f^{-1}(\mathcal{O})$ 是开的.

反之, 假定每个开集的原象是开的. 令 $x_{0}\in\mathbf{R}$ 且 $\varepsilon>0$. 集合 $(f(x_{0})-\varepsilon,f(x_{0})+\varepsilon)$ 是开的, 因此 $f^{-1}((f(x_{0})-\varepsilon,f(x_{0})+\varepsilon))$ 是开的, 且包含 $x_{0}$. 因此存在 $\delta>0$ 使得 $(x_{0}-\delta,x_{0}+\delta)\subseteq f^{-1}((f(x_{0})-\varepsilon,f(x_{0})+\varepsilon))$. 即对 $|x-x_{0}|<\delta$, $|f(x)-f(x_{0})|<\varepsilon$. 所以 $f$ 在 $x_{0}$ 连续.
\end{proof}

\begin{theorem}[极值定理]\label{theorem:ch01-极值}
令 $f$ 为定义在闭有界实数集 $F$ 上的连续实值函数. 则 $f$ 在 $F$ 上取到最大值与最小值.
\end{theorem}

\begin{proof}
令 $M=\sup_{x\in F}f(x)$. 根据 Heine-Borel 定理与 Bolzano-Weierstrass 定理, 存在序列 $\{x_{n}\}\subseteq F$ 使得 $f(x_{n})\to M$. 由于 $F$ 是闭有界的, Bolzano-Weierstrass 定理给出子序列 $x_{n_{k}}\to x_{0}\in F$. 由连续性, $f(x_{n_{k}})\to f(x_{0})$. 但 $f(x_{n_{k}})\to M$, 故 $f(x_{0})=M$. 最小值类似.
\end{proof}

\begin{theorem}[介值定理]\label{theorem:ch01-介值}
令 $f:[a,b]\to\mathbf{R}$ 连续. 若 $f(a)<0$ 且 $f(b)>0$, 则存在 $c\in(a,b)$ 使得 $f(c)=0$.
\end{theorem}

\begin{proof}
定义 $E=\{x\in[a,b]\mid f(x)<0\}$. $E$ 非空且有上界 $b$. 令 $c=\sup E$. 我们宣称 $f(c)=0$. 首先 $c\in[a,b]$. 若 $f(c)<0$, 由连续性, 存在 $\delta>0$ 使得对 $x\in(c-\delta,c+\delta)$, $f(x)<0$. 于是 $c+\delta/2\in E$, 与 $c$ 为上界矛盾. 若 $f(c)>0$, 由连续性, 存在 $\delta>0$ 使得对 $x\in(c-\delta,c+\delta)$, $f(x)>0$. 于是 $c-\delta/2$ 是 $E$ 的上界, 与 $c=\sup E$ 矛盾. 因此 $f(c)=0$.
\end{proof}

\begin{definition}[一致连续]\label{definition:ch01-一致连续}
函数 $f:E\to\mathbf{R}$ 称为在 $E$ 上一致连续, 若对每个 $\varepsilon>0$, 存在 $\delta>0$, 使得对所有 $x,x'\in E$,
\begin{align*}
|x-x'|<\delta\Longrightarrow|f(x)-f(x')|<\varepsilon.
\end{align*}
\end{definition}

\begin{theorem}\label{theorem:ch01-一致连续}
定义在闭有界实数集 $F$ 上的连续函数 $f:F\to\mathbf{R}$ 在 $F$ 上一致连续.
\end{theorem}

\begin{proof}
用反证法. 若 $f$ 在 $F$ 上不一致连续, 则存在 $\varepsilon_{0}>0$, 使得对每个 $\delta>0$, 存在 $x,x'\in F$ 满足 $|x-x'|<\delta$ 但 $|f(x)-f(x')|\geqslant\varepsilon_{0}$. 对每个自然数 $n$, 取 $\delta=1/n$, 选取 $x_{n},x_{n}'\in F$ 满足 $|x_{n}-x_{n}'|<1/n$ 但 $|f(x_{n})-f(x_{n}')|\geqslant\varepsilon_{0}$. 由于 $F$ 是有界的, Bolzano-Weierstrass 定理给出子序列 $x_{n_{k}}\to x_{0}\in F$. 则 $x_{n_{k}}'\to x_{0}$ 也成立, 因为 $|x_{n_{k}}-x_{n_{k}}'|<1/n_{k}\to 0$. 由连续性, $f(x_{n_{k}})\to f(x_{0})$ 且 $f(x_{n_{k}}')\to f(x_{0})$. 因此 $|f(x_{n_{k}})-f(x_{n_{k}}')|\to 0$, 与 $|f(x_{n_{k}})-f(x_{n_{k}}')|\geqslant\varepsilon_{0}$ 矛盾.
\end{proof}

\begin{definition}[递增、递减、单调]\label{definition:ch01-单调函数}
函数 $f:E\to\mathbf{R}$ 称为递增的, 若对所有 $x<x'$（$x,x'\in E$）, $f(x)\leqslant f(x')$. 称为递减的, 若对所有 $x<x'$, $f(x)\geqslant f(x')$. 单调指的是或者递增或者递减.
\end{definition}

对于函数 $f$ 和点 $x_{0}\in\mathbf{R}$, 定义右极限 $f(x_{0}+)$ 与左极限 $f(x_{0}-)$ 为
\begin{align*}
f(x_{0}+)=\lim_{x\downarrow x_{0}}f(x),\qquad f(x_{0}-)=\lim_{x\uparrow x_{0}}f(x),
\end{align*}
若这些极限存在. 若 $f(x_{0}+)$ 与 $f(x_{0}-)$ 都存在但 $f(x_{0}+)\neq f(x_{0}-)$, 称 $f$ 在 $x_{0}$ 有跳跃不连续.

\begin{exercise}\label{exercise:ch01-47}
给出在 $\mathbf{R}$ 上连续但在 $0$ 不是 Lipschitz 的函数的例子.
\end{exercise}

\begin{exercise}\label{exercise:ch01-48}
证明: 若 $f,g:E\to\mathbf{R}$ 都连续, 则 $f+g$, $f-g$, $fg$, $|f|$, $\max\{f,g\}$, $\min\{f,g\}$ 都连续. 若 $g$ 在 $E$ 上不取零, $f/g$ 也连续.
\end{exercise}

\begin{exercise}\label{exercise:ch01-49}
证明: 若 $f:[a,b]\to\mathbf{R}$ 连续且 $f(a)\leqslant c\leqslant f(b)$, 则存在 $x\in[a,b]$ 使得 $f(x)=c$.
\end{exercise}

\begin{exercise}\label{exercise:ch01-50}
证明: 若 $f:[0,1]\to[0,1]$ 连续, 则存在 $x\in[0,1]$ 使得 $f(x)=x$.
\end{exercise}

\begin{exercise}\label{exercise:ch01-51}
给出在 $(0,1)$ 上连续但不是一致连续的函数的例子.
\end{exercise}

\begin{exercise}\label{exercise:ch01-52}
证明: 若 $f:(a,b)\to\mathbf{R}$ 连续且单调递增, 则 $f$ 在 $(a,b)$ 上至多有可数个不连续点.
\end{exercise}

\begin{exercise}\label{exercise:ch01-53}
令 $f$ 为定义在闭有界集 $F$ 上的连续函数. 证明: 对每个 $\varepsilon>0$, 存在定义在 $\mathbf{R}$ 上的连续函数 $g$, 使得对所有 $x\in F$, $|f(x)-g(x)|<\varepsilon$.
\end{exercise}

\begin{exercise}\label{exercise:ch01-54}
证明: 若 $f$ 在 $E$ 上一致连续, 且 $\{x_{n}\}\subseteq E$ 是 Cauchy 序列, 则 $\{f(x_{n})\}$ 也是 Cauchy 序列.
\end{exercise}

\begin{exercise}\label{exercise:ch01-55}
证明: 若 $f$ 在开区间 $(a,b)$ 上一致连续, 则 $f(a+)$ 与 $f(b-)$ 存在.
\end{exercise}

\begin{exercise}\label{exercise:ch01-56}
给出函数 $f:\mathbf{R}\to\mathbf{R}$ 的例子, 使得 $f^{-1}(\mathcal{O})$ 对每个开集 $\mathcal{O}$ 是开的, 但 $f$ 不连续.
\end{exercise}

\begin{exercise}\label{exercise:ch01-57}
证明: 若 $f:\mathbf{R}\to\mathbf{R}$ 连续, 且对每个闭有界区间 $[a,b]$, $f([a,b])$ 是闭的, 则 $f$ 是单调的.
\end{exercise}

\begin{exercise}\label{exercise:ch01-58}
令 $f:\mathbf{R}\to\mathbf{R}$ 连续. 证明: 若对每个开集 $\mathcal{O}$, $f(\mathcal{O})$ 是开的, 则 $f$ 是严格单调的.
\end{exercise}

\begin{exercise}\label{exercise:ch01-59}
证明: 每个单调函数的不连续点集是可数的.
\end{exercise}

\begin{exercise}\label{exercise:ch01-60}
令 $f:[a,b]\to\mathbf{R}$ 连续. 证明: $f$ 是 Lipschitz 的当且仅当 $f$ 在 $(a,b)$ 上有有界导数.
\end{exercise}

\end{document}